%% ============================================================
%%  REFERÊNCIA — pacote handoutAprendizadoRL
%%  Copie este arquivo para começar um handout novo.
%%
%%  Compilar:  latexmk -xelatex referencia-aula03.tex
%% ============================================================
\documentclass[11pt, a4paper]{article}

%% .sty do projeto na raiz do repositório (../)
\makeatletter
\def\input@path{{./}{../}}
\makeatother

\usepackage{handoutAprendizadoRL}
\usepackage{babel}
\babelprovide[import, main]{brazilian}

\setaula{Aula N — Título da folha}
\setprof{Prof. Marcelo Keese Albertini}

%% Macros locais (mesmas da apresentação, quando houver)
\newcommand{\Bop}{\mathcal{B}}   % exemplo: operador de Bellman

\begin{document}

\capa{Aula N: Título do handout}
     {Subtítulo curto \textbullet\ tópicos principais}
     {Prof. Marcelo Keese Albertini}
     {Adaptado de \ldots\ \textbullet\ 2026}

\section{Caixas temáticas}

As mesmas caixas do deck Beamer, com as mesmas assinaturas — o conteúdo
pode ser copiado de um formato para o outro sem edição.

\begin{defbox}{Horizonte $H$}
  Número de passos de tempo em cada episódio.
\end{defbox}

\begin{defbox}{}
  Com título vazio, a caixa escreve apenas ``Definição''.
\end{defbox}

\begin{teobox}{Proposição — melhoria monotônica}
  $V^{\pol_{i+1}} \ge V^{\pol_i}$, com desigualdade estrita se $\pol_i$ é subótima.
\end{teobox}

\begin{algbox}{Iteração de Valor}
  \begin{itemize}
    \item Inicialize $V_0(s) = 0$ para todo $s$.
    \item Repita até convergir:
      $V_{k+1}(s) = \max_a\bigl[\rec(s,a) + \gamma\sum_{s'}\tran{s'}{s,a}V_k(s')\bigr]$.
  \end{itemize}
\end{algbox}

\begin{ideiabox}
  O valor de hoje se decompõe em: o que ganho agora mais o valor de onde posso parar.
\end{ideiabox}

\begin{alertabox}
  Busca exaustiva é inviável: há $\abs{\gA}^{\abs{\gS}}$ políticas determinísticas.
\end{alertabox}

\begin{exbox}{— Mars Rover}
  Calcule $V_{k+1}(s_6)$ com $\gamma = 0{,}9$.
\end{exbox}

\begin{provabox}
  Toda contração em norma $\infty$ converge para um ponto fixo único.
\end{provabox}

\begin{discbox}
  Um fator de desconto $\gamma$ grande dá mais peso a recompensas de curto prazo?
\end{discbox}

\section{Macros matemáticas}

As mesmas do deck:
\begin{itemize}
  \item estados, ações, políticas e valores: $\gS$, $\gA$, $\pol$, $\polstar$,
        $\Vpi$, $\Qpi$, $\Vstar$;
  \item probabilidade e álgebra: $\E$, $\Prob$, $\rec$, $\tran{s'}{s,a}$,
        $\abs{\gS}$, $\norminf{V_k - V_j}$, $\defeq$, $\argmax_a Q(s,a)$.
\end{itemize}

Realce dentro de fórmula e anotação sob um termo:
\[
  \mdestaque{rlLaranja}{V = (I - \gamma P)^{-1}\rec}
  \quad
  V(s) = \termo{\rec(s)}{recompensa imediata}
       + \termo{\gamma\textstyle\sum_{s'}\tran{s'}{s}V(s')}{futuro descontado}
\]

\section{Diagramas}

Os estilos tikz do pacote (\texttt{estado}, \texttt{trans},
\texttt{estado terminal}, \ldots) desenham um MDP em poucas linhas:

\begin{center}
  \begin{tikzpicture}[node distance=12mm]
    \node[estado] (s1) {$s_1$};
    \node[estado, right=of s1] (s2) {$s_2$};
    \node[estado terminal, right=of s2] (s3) {$s_3$};
    \draw[trans] (s1) -- (s2);
    \draw[trans destac] (s2) -- node[above, font=\scriptsize, color=rlLaranja] {$a_1$} (s3);
  \end{tikzpicture}
\end{center}

\section{Tabelas}

Prefira \texttt{booktabs} com larguras fixas:

\begin{center}\footnotesize
\begin{tabular}{@{}lcc@{}}
  \toprule
  & $V(B)$ & $V(A)$\\
  \midrule
  Monte Carlo & $0{,}75$ & $0$\\
  TD(0) em lote & $0{,}75$ & $0{,}75$\\
  \bottomrule
\end{tabular}
\end{center}

\section{Ajustes finos}

\begin{itemize}
  \item Margens: chame \texttt{\textbackslash geometry\{...\}} \emph{depois} de
        carregar o pacote.
  \item Cores pelo nome, com a mistura do xcolor (\texttt{rlAzul!15}):
        \texttt{rlAzul}, \texttt{rlAzulClaro}, \texttt{rlTeal},
        \texttt{rlLaranja}, \texttt{rlAmbar}, \texttt{rlVermelho},
        \texttt{rlVerde}, \texttt{rlCinza}, \texttt{rlFundo},
        \texttt{rlFundoQuiz}.
  \item Toda caixa aceita opções do tcolorbox no argumento opcional, como
        \texttt{[fontupper=\textbackslash small]}.
  \item Macros que a aula precisar e que não existem no pacote entram no
        bloco ``Macros locais'', logo após \texttt{\textbackslash setprof}.
\end{itemize}

\begin{teobox}[fontupper=\small]{Caixa com opção de tcolorbox}
  Este texto está menor porque a caixa recebeu \texttt{fontupper} no argumento
  opcional.
\end{teobox}

\end{document}
